\section{序論}
8月16日朝から17日朝にかけてのXでは、モデルそのものの大型発表よりも、生成AIを取り巻く事業基盤と運用体制の変化が目立った。中でもStripeによるOpenRouter買収報道、OpenAIのPreparednessチーム再編報道、AWSによるエージェント向けポリシー言語Dogwoodの公開は、生成AIの競争軸がモデル性能だけでなく、課金・ルーティング、安全統治、エージェント実行基盤へ広がっていることを印象づけた。

同時に、Claudeのテキスト透かしをめぐる規制対応、AIチャットボットを用いた詐欺研究、ChatGPT Mac版の操作履歴機能に関する懸念など、安全・プライバシー面の話題も継続した。政策面では米国と中国を軸とするAI枠組みの競合が話題となり、研究コミュニティでは長時間エージェント設計やActive Inference向け生成モデルに関する論文共有が続いた。

全体として、この24時間は「新モデルの日」というより、AIをどう配備し、統治し、商流へ組み込むかが前面に出た観測日だった。
\dailyxnote{Grok入力には07:00 JSTの締切直後と記録された投稿も一部含まれる。Daily版ではauthoritative intakeの採用順を維持し、厳密な採否・タイムライン検証は週刊版の工程に委ねる。}
